% ============================================================
% CAPÍTULO 2 — DESCRIPCIÓN GENERAL (IEEE 830 §2)
% ============================================================
\chapter{Descripción general}\label{cap:descripcion}

\section{Perspectiva del producto}

El \nombreSistema{} es un producto autónomo (no forma parte de un sistema
mayor) construido sobre una arquitectura web cliente--servidor de tres capas:

\begin{description}
    \item[Capa de presentación (cliente).] HTML5, CSS3 y JavaScript sin
    frameworks. La navegación sigue el patrón \textit{shell + frames}: existe
    una única página contenedora (\texttt{paginas/dashboard/index.html}) que
    renderiza una sola vez la barra lateral de navegación y la barra
    superior, y carga cada Frame dentro de un marco interno
    (\texttt{iframe}) sin recargar la estructura general. El estilo visual
    (\textit{glassmorphism}: paneles flotantes translúcidos con desenfoque
    sobre una imagen de fondo con capa oscura) se hereda de la pantalla de
    inicio de sesión hacia todo el sistema.

    \item[Capa de lógica (servidor).] PHP~7.3 con acceso a datos mediante
    PDO. Cada operación de negocio se expone como un endpoint independiente
    que recibe y responde JSON, siguiendo el patrón uniforme: verificar
    sesión $\rightarrow$ verificar permiso RBAC $\rightarrow$ verificar
    disponibilidad global del módulo $\rightarrow$ ejecutar consulta
    $\rightarrow$ responder \texttt{\{ok, msg, data\}}. Las bibliotecas
    compartidas son \texttt{servidor/config.php} (conexión, sesión,
    autorización), \texttt{servidor/permisos/registro.php} (catálogo central
    de módulos y acciones) y \texttt{servidor/menu/semilla.php} (menú
    inicial de usuarios nuevos).

    \item[Capa de datos.] Base de datos MySQL~8.0 (\texttt{proyecto\_pw})
    con 11 tablas, integridad referencial completa mediante claves foráneas
    y restricciones de unicidad sobre las claves de negocio.
\end{description}

El estado de sesión del cliente (token, usuario, permisos, menú) se conserva
en \texttt{sessionStorage} y se revalida contra el servidor en cada carga de
página.

\subsection{Componentes del cliente}

\begin{longtable}{|>{\ttfamily}p{4.2cm}|p{10.4cm}|}
\hline
\rowcolor{azulcorp!15} \textbf{Archivo} & \textbf{Responsabilidad} \\ \hline
\endhead
config.js & Constantes globales, rutas de páginas y catálogo de endpoints. \\ \hline
utils.js & Peticiones JSON, alertas flotantes, indicador de carga, formato y
           saneamiento de texto (prevención XSS). \\ \hline
session.js & Lectura/escritura del estado de sesión y verificación de
             permisos en el cliente. \\ \hline
router.js & Guardias de navegación: protección de páginas, verificación de
            permisos y de disponibilidad del módulo, redirecciones. \\ \hline
validaciones.js & Validación de formularios y política de contraseñas. \\ \hline
shell.js & Controlador del shell: carga de Frames en el marco, título
           dinámico, ítem activo, enlace profundo por \textit{hash}. \\ \hline
menu.js & Renderizado de la barra lateral (acordeones) y barra superior. \\ \hline
auth.js & Inicio de sesión, captcha, bloqueo por intentos, mostrar/ocultar
          contraseña, cierre de sesión. \\ \hline
dashboard.js, usuarios.js, roles.js, permisos.js, menu-personal.js,
menu-config.js, cambiar-password.js, frame1.js & Controlador de cada Frame. \\ \hline
\end{longtable}

\section{Funciones del producto}

El sistema ofrece doce pantallas (Frames) organizadas en los siguientes
grupos funcionales:

\begin{longtable}{|p{3.4cm}|p{11.2cm}|}
\hline
\rowcolor{azulcorp!15} \textbf{Grupo} & \textbf{Funciones} \\ \hline
\endhead
Autenticación & Inicio de sesión con captcha y bloqueo por intentos; cierre
                de sesión; expiración y verificación de sesión; cambio de
                contraseña obligatorio en el primer acceso. \\ \hline
Navegación & Shell con menú lateral filtrado por permisos, disponibilidad
             global y organización personal; carga de Frames sin recarga;
             modo móvil. \\ \hline
Dashboard & Saludo contextual, accesos rápidos personalizables por usuario,
            estadísticas del sistema (solo Administrador) y frase motivadora
            aleatoria (solo usuarios estándar). \\ \hline
Administración & Gestión de usuarios (listar, filtrar, crear, editar,
                 activar/desactivar), de roles (listar, crear, editar) y de
                 permisos por rol en dos niveles (acceso a Frames y acciones
                 internas). Exclusiva del Administrador. \\ \hline
Menús & Personalización individual del menú mediante arrastrar y soltar
        (SuperMenus propios) y administración de la disponibilidad global de
        ItemMenus (solo Administrador). \\ \hline
Mi Perfil & Consulta de los datos propios y cambio de contraseña con
            indicador de fortaleza. \\ \hline
Frames de ejemplo & Frame~1: CRUD de tareas con permisos por operación;
                    Frames~2--5: pantallas de demostración con control de
                    acceso. \\ \hline
\end{longtable}

\section{Características de los usuarios}

\begin{longtable}{|p{2.6cm}|p{5.6cm}|p{5.9cm}|}
\hline
\rowcolor{azulcorp!15} \textbf{Actor} & \textbf{Descripción} &
\textbf{Privilegios} \\ \hline
\endhead
Visitante & Persona no autenticada. No requiere conocimientos previos. &
  Únicamente puede acceder a la pantalla de inicio de sesión; cualquier otra
  página lo redirige al login. \\ \hline
Usuario estándar & Usuario autenticado con alguno de los roles Rol1--Rol7.
  Usuario final con conocimientos básicos de navegación web. &
  Accede solo a los Frames habilitados a su rol; siempre dispone de
  Dashboard, Mi Perfil, la lista de usuarios en modo lectura y la
  personalización de su propio menú; dentro de cada Frame solo ve las
  acciones autorizadas. \\ \hline
Administrador & Usuario con el rol \texttt{Admin} (identificador de rol~1).
  Responsable de la operación del sistema. &
  Posee implícitamente todos los permisos (no es configurable); es el único
  que administra usuarios, roles, permisos y la disponibilidad global de los
  menús. \\ \hline
\end{longtable}

\section{Restricciones}

\begin{itemize}
    \item El servidor requiere PHP~7.3 o superior con PDO/MySQL y una base
          de datos MySQL~8.0 o MariaDB equivalente.
    \item El cliente no utiliza frameworks ni bibliotecas externas
          (JavaScript puro, CSS propio); no existen dependencias de
          terceros ni gestores de paquetes.
    \item Toda la comunicación cliente--servidor se realiza por HTTP con
          peticiones POST y cuerpo JSON; el token de sesión viaja en el
          cuerpo de cada petición.
    \item Las contraseñas se almacenan exclusivamente con hash bcrypt
          (costo 12); ninguna operación devuelve la contraseña.
    \item La autorización se valida siempre en el servidor; las
          comprobaciones del cliente son únicamente de experiencia de
          usuario.
    \item Los módulos del sistema deben estar declarados en el registro
          central de Frames; el guardado de permisos rechaza módulos o
          acciones no registrados.
    \item Las rutas del proyecto están ancladas al prefijo
          \texttt{/NRC30713-Web/Proyecto\_pw/} del servidor web.
\end{itemize}

\section{Suposiciones y dependencias}

\begin{itemize}
    \item Se asume un servidor Apache local (entorno de referencia: AppServ
          en el puerto 8080) con la base de datos \texttt{proyecto\_pw}
          creada e inicializada con los datos semilla.
    \item Se asume el uso de un navegador moderno con JavaScript,
          \texttt{sessionStorage} e \texttt{iframe} habilitados.
    \item La personalización del menú mediante arrastrar y soltar utiliza la
          API nativa de HTML5 \textit{Drag \& Drop}, orientada a dispositivos
          con puntero (escritorio).
    \item La existencia del usuario administrador inicial
          (\texttt{admin}, rol~1) es una precondición de operación: el
          sistema no ofrece registro público de cuentas.
    \item Los relojes del servidor determinan la expiración de sesiones y
          bloqueos; se asume una zona horaria consistente.
\end{itemize}
